A 3D nyomtatási folyamat során különböző hibák léphetnek fel, amelyek anyagveszteséget, időveszteséget és akár a berendezés károsodását is okozhatják.
A dolgozat célja egy olyan valós idejű monitorizáló rendszer fejlesztése, amely képfeldolgozási módszerekkel képes felismerni a nyomtatási hibákat.
A javasolt megoldás egy többkamerás megfigyelőrendszerre épül, amely a nyomtatási folyamatot folyamatosan rögzíti, majd a képkockákat egy felhőalapú
környezetben futó YOLOv8 neurális hálózat dolgozza fel. A modell egy saját annotált adathalmazon került betanításra, és 9 hibaosztály felismerésére képes, és 89.2\% mAP50 pontossággal
és 83.6\% recall értékkel rendelkezik, illetve másodpercenként 6 képkockát dolgoz fel 23 ms inferencia idővel.
A rendszer része egy értesítési mechanizmus is, amely anomália észlelése esetén azonnali visszajelzést ad a felhasználónak.
A kísérleti eredmények alapján a rendszer alkalmas valós idejű hibadetektálásra, és hozzájárulhat a nyomtatási folyamat hatékonyságának növeléséhez,
valamint az anyagveszteség csökkentéséhez.\\



\textbf{Kulcsszavak}: 3D nyomtatás, Konvolúciós neuronháló, YOLOv8x, Raspberry Pi, Google Cloud, Vertex AI, hiba, Terraform, Docker, Firebase\\